\section{Record topology and differentiated moving tests}
\label{sec:g-record-response}
A record limit is not, by itself, a response theorem. We specify both the
metric and the differentiable test class. This also fixes the meaning of
the bounded Lipschitz tests at the end of Theorem~\ref{thm:record}.
The unaltered circular estimates in that theorem continue to apply.

In a fixed oriented channel put $s=\sqrt{2d}$ and use the exact Morse
coordinates $w=sz$, together with first residual time $r=d\eta$.
The latent domain is the same $\mathcal U$ as in \eqref{eq:v3-latent-domain}.
Its unweighted conditional density relative to $\mathsf m$ is
\begin{equation}\label{eq:g-latent-density}
 f_{j,e,\xi,s}(z,\eta)=
 \frac{a_{j,e}(\xi,sz)}{\displaystyle\int_{\mathcal U}
                         a_{j,e}(\xi,sz)\dd\mathsf m}.
\end{equation}
This is a density on common latent coordinates, not a total-variation
comparison between physical measures on distinct embedded record surfaces.
In general $f=1+O(s)$; for the circular unweighted density its evenness
gives the stronger $O(s^2)$ already proved in Theorem~\ref{thm:record}.

Let $\mathcal R_{j,e}(\xi,z,\eta,s)$ be the array consisting of
$y_i/s$, $(t_i-ig_e)/d$, and both transverse velocity traces divided
by $s$, for $0\le i\le j$. The transverse frame is that of the closest
chord; the axial values have zero transverse velocity. At $s=0$ use
the continuous limits. Equip each such array with the maximum norm,
\begin{equation}\label{eq:g-record-metric}
 d_{\infty,j}(\mathcal R,\widetilde{\mathcal R})
 =\max_{0\le i\le j}
    \max\{|Y_i-\widetilde Y_i|,|T_i-\widetilde T_i|,
          |V_i^--\widetilde V_i^-|,|V_i^+-\widetilde V_i^+|\}.
\end{equation}
For tests involving only positions and times omit the other entries.
An unnormalized Euclidean norm of the entire growing array is not used.
The limiting positions are the linear bridge \eqref{eq:g-bridge} applied
to $\mathsf H_{j,e}^{-1/2}z$. Its limiting excess flight times are
\[
 Q_i=\kappa_{i\bmod2}Y_i^2+\kappa_{1-(i\bmod2)}Y_{i+1}^2
                          +(Y_{i+1}-Y_i)^2/g_e,
 \qquad \sum_{i<j}Q_i=|z|^2.
\]
Thus the limiting time array is $\eta+\sum_{h<i}Q_h$.

\begin{proposition}[Differentiated smooth record tests]\label{prop:g-smooth-tests}
Fix any derivative order $m$. Let $\varphi_{j,\xi}$ be smooth tests of
these scaled arrays with bounded derivatives through order $m+1$, in
operator norm for the maximum-norm array space, and with the corresponding
mixed parameter derivatives bounded independently of $j$. Then the
conditional expectations $I_j(\xi,s)$ are smooth in $(\xi,s)$ at $s=0$
on the right, and
\begin{equation}\label{eq:g-differentiated-record}
 \sup_{\xi,e,j}\bigl|\partial_\xi^\alpha
                  (I_j(\xi,s)-I_j(\xi,0))\bigr|\le C_m s,
 \qquad |\alpha|\le m.
\end{equation}
A sufficient test class consists of a fixed finite number of selected
record entries and smooth functions with uniformly bounded derivatives.
For uniformly bounded Lipschitz tests without differentiability, the
same statement without $\partial_\xi^\alpha$ follows from
\eqref{eq:g-record-metric}. The circular unweighted convergence in
Theorem~\ref{thm:record} retains its stronger $O(d)$ rate for its
positions and times.
\end{proposition}
\begin{proof}
The bridge, Morse inverse, chord directions and reflections are smooth
in $\xi,w$. They vanish to the required order at $w=0$: positions and
transverse velocities vanish to first order, and every excess flight
length to second order. Taylor's integral formula therefore extends
the divisions by $s$ and $s^2/2=d$ smoothly to $s=0$.
For times, the estimates are summed along the bridge. The sum of the
endpoint-decaying bounds from Lemma~\ref{lem:g-relative} is independent
of $j$. Hence $\mathcal R_j$, all its fixed derivatives, and
$s^{-1}(\mathcal R_j(s)-\mathcal R_j(0))$ have uniform maximum-norm
bounds. The same claims hold for the conditional density
\eqref{eq:g-latent-density}; its denominator is bounded below.

Now write the expectation exactly on the fixed domain:
\[
 I_j(\xi,s)=\int_{\mathcal U}
 f_{j,e,\xi,s}(z,\eta)\,
 \varphi_{j,\xi}(\mathcal R_{j,e}(\xi,z,\eta,s))\dd\mathsf m.
\]
The chain rule in the maximum norm has no dimension factor by the
hypothesis on the test derivatives. Dominated differentiation and the
fundamental theorem of calculus in $s$ prove
\eqref{eq:g-differentiated-record}. This differentiates the actual
parametrization and density; it does not differentiate a total-variation
inequality. For Lipschitz tests, instead bound the difference by the
pathwise maximum-norm error plus the density error. In the circular
positions--times specialization these errors are $O(d)$ as already
proved. The entire statement also holds under bounded smooth additive
sources after replacing the density by its positive normalized factor
$e^{qV_{j,e}}$, with the same general $O(s)$ rate.
\end{proof}

For a physical sampling cut, write the exact excess time of impact $k$
as $T_{j,k}(\xi,z,s)=\sum_{i<k}(\ell_i-g_e)/d$.
It has all fixed derivative bounds uniformly in $j,k$. A sample at
$t^{\rm s}=kg_e+d\,a(\xi)$ precedes that impact precisely when
\[
                 \eta>\psi_{j,k}(\xi,z,s):=
                               a(\xi)-T_{j,k}(\xi,z,s).
\]
Here $a$ is smooth with bounded derivatives; the physical sample moves
with the geometry, and is not an extra random coordinate.

\begin{proposition}[A differentiated transverse sampling interface]
\label{prop:g-moving-cut}
Let a smooth test be supported in a fixed compact subset of $|z|<1$,
and suppose that on this support, for all parameters under consideration,
\begin{equation}\label{eq:g-transversality}
       \nu\le\psi_{j,k}(\xi,z,s)\le1-|z|^2-\nu
\end{equation}
for one fixed $\nu>0$. For a smooth test as in
Proposition~\ref{prop:g-smooth-tests}, its conditional expectation
restricted to the before-impact event is smooth in $\xi,s$, with
uniform fixed-order bounds. If $F_j(\xi,z,\eta,s)$ denotes the
conditional density times the test, relative to $\dd z\dd\eta$, then
its first derivative in any scalar parameter direction is exactly
\begin{equation}\label{eq:g-moving-formula}
 \partial_\xi I_j=
 \int\left\{\int_{\psi_{j,k}}^{1-|z|^2}\partial_\xi F_j\dd\eta
       -F_j(\xi,z,\psi_{j,k},s)\partial_\xi\psi_{j,k}\right\}\dd z.
\end{equation}
Repeated differentiation gives every fixed higher derivative and the
analogue of \eqref{eq:g-differentiated-record}. The boundary term is
part of the formula.
\end{proposition}
\begin{proof}
The derivative of the cut function with respect to $\eta$ is one.
Condition \eqref{eq:g-transversality} separates it uniformly from the
two other boundaries of the latent domain on the support of the test.
Its integral is therefore exactly
$\int\int_{\psi_{j,k}}^{1-|z|^2}F_j\dd\eta\dd z$.
Leibniz' rule gives \eqref{eq:g-moving-formula}, with no other boundary
piece. Smooth bounded density and test derivatives and the uniform
bounds for $T_{j,k}$ justify repeated differentiation under the outer
integral. The fundamental theorem in $s$ gives its differentiated
limit estimate. This proves response of the selected moving physical
interface, rather than response inferred from its small mismatch strip.
\end{proof}

The margin is nonempty uniformly even for interior impacts. Take $a=1/2$
and support the test on $|z|<1/2$. Every flight excess is nonnegative,
and the exact Morse identity gives $0\le T_{j,k}\le|z|^2$ for every
$k$. Both margins are therefore at least $1/4$. A finite collection of cuts with
a fixed strict order and the same margins is handled by the same interval
integration. At a fixed physical sample time $t^{\rm s}$, instead,
$a=(t^{\rm s}-kg_e)/d$ and its derivatives contain factors $k/d$.
Formula \eqref{eq:g-moving-formula} still holds where the margins hold,
but its order-$m$ bounds can grow as $C_m(1+k/d)^m$. No uniform onset
response for all fixed schedules is asserted. When margins fail, or
for an arbitrary bounded decoder, the convergence theorem alone gives
no derivative estimate. The original finite-cut convergence and the
smooth additive response remain separate valid statements.
